\documentclass[10pt,a4paper]{article}
\usepackage{fontspec}
\usepackage[margin=1.6cm,top=1.4cm,bottom=1.4cm]{geometry}
\usepackage{booktabs,tabularx,xcolor}
\pagestyle{empty}
\setlength{\parindent}{0pt}
\setlength{\parskip}{4pt}
\definecolor{ok}{HTML}{1B7F3B}
\definecolor{ko}{HTML}{B3261E}
\begin{document}

{\Large\bfseries Test against a real group of the 2022 challenge: \texttt{smokeonthewatermark}}\\[2pt]
{\small Group \texttt{fakebitinvestigation} — 21/10/2026 — all the numbers are measured}

\textbf{Their scheme.} Haar wavelet with 2 levels (subbands LH2, HL2, HH2), spread spectrum with a local energy mask, LDPC error-correcting code; fixed strength (not calibrated on the WPSNR). Detection: decoded 0/1 bits, cosine with threshold 0.5377. \textbf{Bench:} 12 images (the 3 of the 2022 challenge + 0000, 0010, \dots, 0080), watermarked with their code and their watermark (WPSNR 56.23–58.16 dB, mean 57.60) and with ours (62.05 dB).

\medskip
\textbf{Run 1 — our attack script against them} (the automatic search \texttt{attack\_image()} of \texttt{attacks\_fakebitinvestigation.py}, at the time in a separate script; challenge settings: 25 families + 7 zones + block greedy with their detection as oracle; 12 processes, 29.4 min per image)

\begin{tabularx}{\linewidth}{@{}l r X@{}}
\toprule
image & WPSNR (dB) & lightest successful attack \\
\midrule
0070 & 47.79 & median 7$\times$7, only in the chosen zones \\
0010 & 41.30 & blur $\sigma$ 1 + JPEG QF 16, zones \\
0050 & 41.17 & resize 0.263 (cubic), zones \\
0040 & 39.98 & resize 0.325 (cubic), zones \\
0080 & 39.85 & blur $\sigma$ 0.8 + JPEG QF 19 + blur $\sigma$ 0.6, zones \\
0030 & 39.11 & blur $\sigma$ 1 + JPEG QF 19, zones \\
0060 & 38.95 & resize 0.5 (lab) + JPEG QF 13, zones \\
tree & 38.84 & resize 0.375 (cubic), zones \\
0020 & 38.49 & JPEG QF 8 + blur $\sigma$ 0.8, zones \\
rollercoaster & 37.82 & median 5$\times$5 + sharpening $\sigma$ 4 $\alpha$ 0.3, zones \\
0000, buildings & — & \textcolor{ko}{not broken} (families exhausted in $\sim$9 min, $\sim$190 calls) \\
\bottomrule
\end{tabularx}

\textcolor{ok}{\textbf{Broken 10/12}}, mean WPSNR \textbf{40.33 dB} (median 39.48; min 37.82; max 47.79). Total time 19.9 min, 0 errors, 463–606 calls per broken image. \emph{Independent check:} every saved image, passed again to their detection, gives 0, and the WPSNR recomputed with the lab code matches the one in the log to the hundredth.

\medskip
\textbf{Run 2 — their attacks against our images} (the identical functions of their \texttt{attack.py}; their tool is manual, so we emulated it generously: a grid of 154 attacks and chains — the parameters suggested in their comments, widened, plus blur$\to$JPEG as in their example — and for every image the successful one with the highest WPSNR according to our detection, $\tau$ = 0.087)

\begin{tabularx}{\linewidth}{@{}l r X r@{}}
\toprule
image & WPSNR (dB) & lightest successful attack & successful / 154 \\
\midrule
0070 & 38.74 & JPEG QF 5 + blur $\sigma$ 0.8 & 12 \\
0040 & 36.58 & JPEG QF 7 & 1 \\
0050 & 36.24 & JPEG QF 5 + blur $\sigma$ 0.8 & 10 \\
0010 & 36.09 & AWGN $\sigma$ 0.01 + JPEG QF 5 & 3 \\
0060 & 35.95 & JPEG QF 6 & 3 \\
0000, 0020, 0030, 0080, buildings, rollercoaster, tree & — & \textcolor{ok}{not broken} & 0 \\
\bottomrule
\end{tabularx}

\textcolor{ko}{\textbf{Broken 5/12}}, mean WPSNR \textbf{36.72 dB} (min 35.95; max 38.74): all with JPEG at quality 5–7, our weak point we already knew about.

\medskip
\textbf{Comparison and honest reading}

\begin{tabularx}{\linewidth}{@{}X r r@{}}
\toprule
 & them, attacked by us & us, attacked with their attacks \\
\midrule
WPSNR of the embedding (mean) & 57.60 dB & 62.05 dB \\
broken images & 10/12 & 5/12 \\
mean WPSNR of the successful attacks & 40.33 dB & 36.72 dB \\
corresponding robustness band & 5 points (38–41) & 6 points (35–38) \\
\bottomrule
\end{tabularx}

\begin{itemize}\itemsep1pt
\item \textbf{Our script works on a scheme different from ours}: 10 images out of 12, almost all above 38 dB. The zone greedy matters: the 10 best attacks are all localized.
\item \textbf{The two runs are not symmetric.} Run 1 is our complete engine (search + zones + oracle, 30 min per image); run 2 is only a grid of global attacks with their functions, without zones. One of their operators choosing the zones by hand (\texttt{img\_slice}) could do better: against our own engine, in the simulation of 19/10, we were broken in 18 images out of 27 at 36.5 dB on average.
\item \textbf{At different WPSNRs.} They have $\sim$4.5 dB more energy in the embedding (3 quality points against 5) and are still easier to break at high WPSNR.
\item Not measured: the false positives of their detection (threshold close to the chance value, 0.5).
\end{itemize}

\end{document}
